\documentclass[11pt]{article}
\usepackage[margin=1in]{geometry}
\usepackage{amsmath,amssymb,amsthm,mathtools}
\usepackage[hidelinks]{hyperref}
\usepackage{microtype}

\newtheorem{theorem}{Theorem}
\newtheorem{corollary}[theorem]{Corollary}
\newtheorem{definition}[theorem]{Definition}
\newtheorem{remark}[theorem]{Remark}
\newcommand{\D}{\mathbb D}
\newcommand{\C}{\mathbb C}
\newcommand{\E}{\mathbb E}

\title{A Source-Dependent Fixed-Volume Milestone for Constructive
Four-Dimensional Yang--Mills}
\author{Working draft}
\date{September 2026}

\begin{document}
\maketitle

\begin{abstract}
We formulate a deliberately limited first observable theorem beyond the known
ultraviolet-stability estimates for lattice Yang--Mills theory.  The target is
cutoff convergence, at fixed physical volume and observation scale, of bounded
gauge-invariant cylinder observables of the terminal gauge-covariant block field.
We separate the genuinely open RG estimate from an elementary but essential
analytic reduction: summable complex-source bounds imply convergence and
factorial estimates for every fixed joint cumulant.  No continuum construction,
point-local field, or mass gap is claimed.
\end{abstract}

\section{Scope}

The ultimate Yang--Mills problem requires local fields corresponding to
gauge-invariant curvature polynomials, the Osterwalder--Schrader axioms,
asymptotically-free short-distance behavior, and a positive gap
\cite{JaffeWitten2006,OsterwalderSchrader1975}.  Reflection positivity of the
Wilson lattice action is an indispensable finite-cutoff input
\cite{OsterwalderSeiler1978}, but it does not establish convergence of a
continuum state.

Balaban's gauge-covariant RG establishes ultraviolet stability through a
small-field analysis, cluster expansions, and a large-field renormalization
operation
\cite{Balaban1987RGI,Balaban1988RGII,Balaban1988Convergent,Balaban1989LargeII}.
The present project asks for a comparison estimate between successive ultraviolet
depths for normalized observable-generating functions.  Such an estimate is not
asserted in the cited stability results.

\section{Fixed-volume setup}

Let (G=SU(2)) and let \(\Lambda_{a,R}=a\mathbb Z^4/R\mathbb Z^4\) approximate a
fixed physical four-torus.  The Wilson measure is
\[
 d\mu_{a,R}(U)=Z_{a,R}^{-1}e^{-S_{a,R}(U)}\prod_{e\in E^+}dU_e,
 \qquad
 S_{a,R}(U)=\beta(a)\sum_p\left(1-\frac12\operatorname{ReTr}U_p\right).
\]
Fix a physical block scale \(\ell\), set \(a_n=\ell L^{-n}\), and denote by
\(Q_n\) the (n)-step gauge-covariant averaging map.  The exact identification
of this notation with Balaban's averaging operation
\cite{Balaban1985Averaging} is an explicit proof obligation.

For coarse closed paths \(C_1,\ldots,C_m\) and a function (h) holomorphic near
\([-2,2]^m\), consider
\[
 \mathcal O_{n,h}(U)=h\bigl(
 \operatorname{TrHol}_{C_1}(Q_nU),\ldots,
 \operatorname{TrHol}_{C_m}(Q_nU)\bigr).
\]
This is bounded and gauge invariant on real (SU(2)) configurations.  It is a
fixed-scale cylinder observable, not a point-local field.

\section{Analytic reduction}

For finitely many such observables let
\[
 Z_n(z)=\E_{\mu_{a_n,R}}\exp\left(\sum_{j=1}^qz_j\mathcal O_{n,j}\right),
 \qquad W_n(z)=\log Z_n(z).
\]
For (R=(R_1,\ldots,R_q)>0\), set
\(\D_R=\{z\in\C^q:|z_j|<R_j\}\).

\begin{theorem}[Summable holomorphic increments]\label{thm:analytic}
Suppose every (W_n\) is holomorphic on \(\D_R\) and
\[
 \lVert W_{n+1}-W_n\rVert_{\D_R}\le\delta_n,
 \qquad \sum_{n\ge0}\delta_n<\infty.
\]
Then (W_n\) converges uniformly to a holomorphic (W_\infty\).  For every
multi-index \(\alpha\),
\[
 \left|\partial^\alpha(W_\infty-W_n)(0)\right|
 \le \alpha!\prod_{j=1}^qR_j^{-\alpha_j}
       \sum_{k=n}^{\infty}\delta_k.
\]
Consequently every fixed joint cumulant converges.
\end{theorem}

\begin{proof}
The increment estimate makes (W_n\) uniformly Cauchy, and its limit is locally
uniform and therefore holomorphic.  Apply the multivariable Cauchy formula to
\(W_m-W_n\) on a smaller concentric polytorus and use
\(\lVert W_m-W_n\rVert\le\sum_{k=n}^{m-1}\delta_k\).  Let (m\to\infty\), then
let the smaller polyradius increase coordinatewise to (R\).
\end{proof}

\begin{theorem}[A sufficient zero-free criterion]\label{thm:log}
Suppose (Z_n(0)=1\) and, for one \(0<\eta<1\),
\(\lVert Z_n-1\rVert_{\D_R}\le\eta\) for every (n\).  Then the principal branch
\(W_n=\log Z_n\) is holomorphic on \(\D_R\), and
\[
 \lVert W_n-W_m\rVert_{\D_R}
 \le(1-\eta)^{-1}\lVert Z_n-Z_m\rVert_{\D_R}.
\]
\end{theorem}

\begin{proof}
The range lies in the convex disc (B(1,\eta)\), on which the principal logarithm
is holomorphic.  Integrating (1/w\) along the segment from (Z_m(z)\) to
(Z_n(z)\) gives the estimate because the segment stays at distance at least
\(1-\eta\) from zero.
\end{proof}

\begin{remark}
Real-source convergence is insufficient.  The functions
\(f_n(z)=\sin(nz)/n\) converge uniformly to zero on bounded real intervals while
\(f_n'(0)=1\).  Similarly, (Z_n(z)=1+nz\) shows why normalization at zero does
not supply a cutoff-independent logarithmic domain.
\end{remark}

\section{Open estimate}

The principal target is a cutoff-independent polydisc and a summable estimate
\[
 \sup_{z\in\D_R}|W_{n+1}(z)-W_n(z)|\le\delta_n,
 \qquad \sum_n\delta_n<\infty,
\]
including all large-field polymers.  Theorem~\ref{thm:analytic} then completes the
source-to-cumulant step.  A partition-function stability bound at each depth does
not imply this comparison estimate.

Positive gradient-flow time offers an alternative bounded observable family, but
the resulting fields are nonlocal over a four-dimensional radius of order
\(\sqrt{8t}\) \cite{Luscher2010,RamosSint2016}.  They are therefore postponed
until the RG-adapted cylinder theorem is understood.

\bibliographystyle{plain}
\bibliography{references}
\end{document}
